\section{Introduction}

As machine learning increasingly affects decisions in domains
protected by anti-discrimination law, there is much interest in
algorithmically measuring and ensuring fairness in machine learning.
In domains such as advertising, credit, employment, education, and criminal
justice, machine learning could help obtain more accurate predictions,
but its effect on existing biases is not well understood. Although reliance on
data and quantitative measures can help quantify and eliminate existing biases,
some scholars caution that algorithms can also introduce new biases or
perpetuate existing ones~\cite{BarocasS16}. In May 2014, the Obama
Administration's Big Data Working Group released a report~\cite{Podesta14}
arguing that discrimination can sometimes ``be the inadvertent outcome of the
way big data technologies are structured and used'' and pointed toward ``the
potential of encoding discrimination in automated decisions''. A subsequent
White House report~\cite{WhiteHouse16} calls for ``equal opportunity by design''
as a guiding principle in domains such as credit scoring.

Despite the demand, a vetted methodology for avoiding discrimination against
\emph{protected attributes} in machine learning is lacking. A na\"ive approach might
require that the algorithm should ignore all protected attributes such as
race, color, religion, gender, disability, or family status. However, this idea
of ``fairness through unawareness'' is ineffective due to the existence of
\emph{redundant encodings}, ways of predicting protected attributes from other
features~\cite{PedreshiRT08}.

Another common conception of non-discrimination is \emph{demographic parity}.
Demographic parity requires that a decision---such as accepting or denying a
loan application---be independent of the protected attribute. In the case of a
binary decision~$\Yhat\in\{0,1\}$ and a binary protected
attribute~$A\in\{0,1\},$ this constraint can be formalized by asking that
$\Pr\{\Yhat=1\mid A=0\}=\Pr\{\Yhat=1\mid A=1\}.$ In other words, membership in a
protected class should have no correlation with the decision. Through its
various equivalent formalizations this idea appears in numerous papers.
Unfortunately, as was already argued by Dwork et al.~\cite{DworkHPRZ12}, the
notion is seriously flawed on two counts. First, it doesn't ensure fairness.
Indeed, the notion permits that we accept qualified applicants in the
demographic~$A=0$, but unqualified individuals in~$A=1,$ so long as the
percentages of acceptance match. This behavior can arise naturally, when there
is little or no training data available within $A=1.$ Second, demographic parity
often cripples the utility that we might hope to achieve. Just imagine the
common scenario in which the target variable~$Y$---whether an individual
actually defaults or not---is correlated with~$A.$ Demographic parity would not
allow the ideal predictor $\Yhat = Y,$ which can hardly be considered
discriminatory as it represents the actual outcome. As a result, the loss in
utility of introducing demographic parity can be substantial.

In this paper, we consider non-discrimination from the perspective of supervised
learning, where the goal is to predict a true outcome~$Y$ from features $X$
based on labeled training data, while ensuring they are ``non-discriminatory''
with respect to a specified protected attribute $A$. As in the usual supervised
learning setting, we assume that we have access to labeled training data, in our
case indicating also the protected attribute $A$. That is, to samples from the
joint distribution of $(X,A,Y)$. This data is used to construct a predictor
$\hat{Y}(X)$ or $\hat{Y}(X,A)$, and we also use such data to test whether it is
unfairly discriminatory.

Unlike demographic parity, our notion always allows for the perfectly accurate
solution of~$\Yhat = Y.$ More broadly, our criterion is easier to achieve the more
accurate the predictor $\Yhat$ is, aligning fairness with the central goal in
supervised learning of building more accurate predictors.

The notion we propose is ``oblivious'', in that it is based only on
the joint distribution, or joint statistics, of the true target $Y$,
the predictions $\hat{Y}$, and the protected attribute $A$.  In
particular, it does not evaluate the features in $X$ nor the
functional form of the predictor $\hat{Y}(X)$ nor how it was derived.
This matches other tests recently proposed and conducted, including
demographic parity and different analyses of common risk scores. In
many cases, only oblivious analysis is possible as the functional form
of the score and underlying training data are not public. The only
information about the score is the score itself, which can then be
correlated with the target and protected attribute.  Furthermore, even
if the features or the functional form are available, going beyond
oblivious analysis essentially requires subjective interpretation or
casual assumptions about specific features, which we aim to avoid.

\subsection{Summary of our contributions}
We propose a simple, interpretable, and actionable framework for measuring and
removing discrimination based on protected attributes. We argue that, unlike
demographic parity, our framework provides a meaningful measure of
discrimination, while demonstrating in theory and experiment that we also
achieve much higher utility.  Our key contributions are as follows:
%
\begin{itemize}
\item We propose an easily checkable and interpretable notion of
  avoiding discrimination based on protected attributes. Our notion
  enjoys a natural interpretation in terms of graphical dependency
  models.  It can also be viewed as shifting the burden of uncertainty
  in classification from the protected class to the decision maker. In
  doing so, our notion helps to incentivize the collection of better
  features, that depend more directly on the target rather then the
  protected attribute, and of data that allows better prediction for
  all protected classes.
%
\item
We give a simple and effective framework for constructing classifiers
satisfying our criterion from an arbitrary learned predictor. Rather than
changing a possibly complex training pipeline, the result follows via a simple
post-processing step that minimizes the loss in utility.
%
\item We show that the Bayes optimal non-discriminating (according to
  our definition) classifier is the classifier derived from any Bayes
  optimal (not necessarily non-discriminating) regressor using our
  post-processing step.  Moreover, we quantify the loss that follows
  from imposing our non-discrimination condition in case the score we
  start from deviates from Bayesian optimality. This result helps to
  justify the approach of deriving a fair classifier via
  post-processing rather than changing the original training process.
%
\item
We capture the inherent limitations of our approach, as well as any other
oblivious approach, through a non-identifiability result showing that different
dependency structures with possibly different intuitive notions of fairness
cannot be separated based on any oblivious notion or test.
\end{itemize}

Throughout our work, we assume a source distribution over $(Y,X,A)$, where $Y$
is the target or true outcome (e.g. ``default on loan''), $X$ are the available
features, and $A$ is the protected attribute. Generally, the features $X$ may be
an arbitrary vector or an abstract object, such as an image. Our work does not
refer to the particular form~$X$ has.

The objective of supervised learning is to construct a (possibly randomized)
predictor $\Yhat=f(X, A)$ that predicts~$Y$ as is typically measured through a
loss function. Furthermore, we would like to require that $\Yhat$ {\em does not
discriminate with respect to $A$}, and the goal of this paper is to formalize
this notion.


\section{Equalized odds and equal opportunity}
%
We now formally introduce our first criterion.
%
\begin{definition}[Equalized odds]
We say that a predictor $\Yhat$ satisfies \emph{equalized odds} with
respect to protected attribute~$A$ and outcome~$Y,$ if
$\Yhat$ and $A$ are independent conditional on~$Y.$
\end{definition}
%
Unlike demographic parity, equalized odds allows $\Yhat$ to depend
on~$A$ but only through the target variable~$Y.$ As such, the
definition encourages the use of features that allow to directly
predict~$Y,$ but prohibits abusing~$A$ as a proxy for~$Y.$

As stated, equalized odds applies to targets and protected attributes taking
values in any space, including binary, multi-class, continuous or structured
settings.  The case of binary random variables $Y,\Yhat$ and~$A$ is of central
importance in many applications, encompassing the main conceptual and technical
challenges. As a result, we focus most of our attention on this case, in which
case equalized odds are equivalent to:
\[
\Pr\Set{\Yhat=1\mid A=0, Y=y} = \Pr\Set{\Yhat=1\mid A=1, Y=y}
,\quad y\in\{0,1\}
\]
For the outcome $y=1,$ the constraint requires that $\Yhat$ has equal \emph{true
  positive rates} across the two demographics $A=0$ and $A=1.$ For
$y=0,$ the constraint equalizes \emph{false positive rates}.  The
definition aligns nicely with the central goal of building highly
accurate classifiers, since $\Yhat=Y$ is always an acceptable
solution. However, equalized odds enforces that the accuracy is
equally high in all demographics, punishing models that perform
well only on the majority.

\subsection{Equal opportunity}
%
In the binary case, we often think of the outcome $Y=1$ as the ``advantaged''
outcome, such as ``not defaulting on a loan'', ``admission to a college'' or
``receiving a promotion''. A possible relaxation of equalized odds 
is to require non-discrimination only within the ``advantaged'' outcome group.
That is, to require that people who pay back their loan, have an equal
opportunity of getting the loan in the first place (without specifying any
requirement for
those that will ultimately default). This leads to a relaxation of our notion
that we call ``equal opportunity''.
%
\begin{definition}[Equal opportunity]
We say that a binary predictor~$\Yhat$ satisfies \emph{equal opportunity}
with respect to $A$ and $Y$ if
$\Pr\set{\Yhat=1\mid A=0, Y=1} = \Pr\set{\Yhat=1\mid A=1, Y=1}\,.$
\end{definition}
%
Equal opportunity is a weaker, though still interesting, notion of
non-discrimination, and thus typically allows for stronger utility as we shall
see in our case study.

\subsection{Real-valued scores}

Even if the target is binary, a real-valued predictive score
$R=f(X,A)$ is often used (e.g. FICO scores for predicting loan
default), with the interpretation that higher values of $R$ correspond
to greater likelihood of $Y=1$ and thus a bias toward predicting
$\Yhat=1$.  A binary classifier $\Yhat$ can be obtained by
thresholding the score, i.e. setting $\Yhat=\mathbb{I}\{R>t\}$ for
some threshold $t$.  Varying this threshold changes the trade-off
between sensitivity (true positive rate) and specificity (true negative rate).

Our definition for equalized odds can be applied also to such score
functions: a score $R$ satisfies equalized odds if $R$ is independent
of $A$ given $Y$.  If a score obeys equalized odds, then any
thresholding $\Yhat=\mathbb{I}\{R>t\}$ of it also obeys equalized odds
(as does any other predictor derived from $R$ alone).  

In Section~\ref{sec:achieving}, we will consider scores that might not satisfy
equalized odds, and see how equalized odds predictors can be derived
from them and the protected attribute $A$, by using different
(possibly randomized) thresholds depending on the value of $A$. 
The same is possible for equality of opportunity without the need for 
randomized thresholds.

\subsection{Oblivious measures}
As stated before, our notions of non-discrimination are \emph{oblivious} 
in the following formal sense.
%
\begin{definition}
  A property of a predictor $\Yhat$ or score $R$ is said to be
  \emph{oblivious} if it only depends on the joint distribution of
  $(Y,A,\Yhat)$ or $(Y,A,R)$, respectively.
\end{definition}
%
As a consequence of being oblivious, all the information we need to verify our
definitions is contained in the \emph{joint distribution} of predictor,
protected group and outcome, $(\Yhat, A, Y).$ In the binary case, when $A$ and
$Y$ are reasonably well balanced, the joint distribution of $(\Yhat, A, Y)$ is
determined by $8$ parameters that can be estimated to very high accuracy from
samples. We will therefore ignore the effect of finite sample perturbations and
instead assume that we know the joint distribution of $(\Yhat, A, Y).$


\section{Comparison with related work}

There is much work on this topic in the social sciences and legal scholarship;
we point the reader to Barocas and Selbst~\cite{BarocasS16} for an excellent
entry point to this rich literature. See also the survey by Romei and
Ruggieri~\cite{RomeiR14}, and the references at
\url{http://www.fatml.org/resources.html}.

In its various equivalent notions, demographic parity
appears in many papers, such
as~\cite{CaldersKP09,Zliobaite15,Zafar15} to name a few.
%
Zemel et al.~\cite{ZemelWSPD13} propose an interesting way of achieving
demographic parity by aiming to learn a representation of the data that is
independent of the protected attribute, while retaining as much information
about the features~$X$ as possible. Louizos et al.~\cite{LouizosSLWZ15} extend
on this approach with deep variational auto-encoders.
%
Feldman et al.~\cite{FeldmanFMSV15} propose a formalization of ``limiting
disparate impact''. For binary classifiers, the condition states that
$\Pr\set{ \Yhat = 1 \mid A = 0} \le 0.8\cdot \Pr\set{\Yhat = 1 \mid A = 1}.$
The authors argue that this corresponds to the ``80\% rule'' in the legal
literature.  The notion differs from demographic parity mainly in that it
compares the probabilities as a ratio rather than additively, and in that it
allows a one-sided violation of the constraint.

While simple and seemingly intuitive, demographic parity has serious conceptual
limitations as a fairness notion, many of which were pointed out in work of
Dwork et al.~\cite{DworkHPRZ12}. In our experiments, we will see that
demographic parity also falls short on utility.  Dwork et al.~\cite{DworkHPRZ12}
argue that a sound notion of fairness must be \emph{task-specific}, and
formalize fairness based on a hypothetical similarity measure $d(x,x')$
requiring similar individuals to receive a similar distribution over outcomes.
In practice, however, in can be difficult to come up with a suitable metric.
Our notion is task-specific in the sense that it makes critical use of the final
outcome~$Y,$ while avoiding the difficulty of dealing with the features~$X.$

In a recent concurrent work, Kleinberg, Mullainathan and Raghavan~\cite{KleinbergMR16} showed that in general a score that is \emph{calibrated within each group} does \emph{not} satisfy a criterion equivalent to equalized odds for binary predictors. This result highlights that calibration alone does not imply non-discrimination according to our measure. Conversely, achieving equalized odds may in general compromise other desirable properties of a score.

Early work of Pedreshi et al.~\cite{PedreshiRT08} and several follow-up works
explore a logical rule-based approach to non-discrimination. These approaches
don't easily relate to our statistical approach.
